% SPDX-License-Identifier: Apache-2.0
% covers: Uart
\datasheet{Uart}{Vreteno's serial port: an AXI-Lite peripheral with SiFive's register map, a line each way, a transmit queue of eight and a receive queue of sixty-four.}

\dsfacts{\code{cpu/vreteno/src/uart.rs}}{\code{vreteno32::uart::Uart}}%
{\code{//docs:vreteno}, \code{//docs:noc}}

\subsection*{Function}

\code{Uart} sends bytes on \code{tx} and receives them on \code{rx}.
A frame is one start bit, eight data bits least significant first and
one stop bit, or two with \code{nstop}. Each bit lasts \code{div} plus
one cycles. The lines rest high.

Its registers are SiFive's \code{sifive,uart0} (issue 1011), so
Linux's \code{serial/sifive.c}, OpenSBI's console and Zephyr's
\code{uart\_sifive} drive it unchanged. A byte written to
\code{txdata} joins a transmit queue of eight. A byte received lands
in a receive queue of sixty-four, and a read of \code{rxdata} takes
the oldest. SiFive's queue holds eight, but Linux's driver on this core
takes longer to answer the interrupt than eight characters take to
arrive at 115200 baud, so a line pasted at its shell lost all but its
first eight (issue 1153). \code{irq} is high while a watermark that \code{ie} enables is
passed.

It is an AXI-Lite peripheral. It sits behind a \code{LiteBridge<1, ..>},
which hands it one transaction at a time with no identifier and no
burst. Its link is one port, \code{bus}, a \code{LitePort} of the
five channels: reads come on \code{bus\_ar} and are answered on
\code{bus\_r}; writes come on \code{bus\_aw} and \code{bus\_w} and
are answered on \code{bus\_b}. The input \code{rst} stops both sides,
empties both queues and puts the controls back to their reset values.

\subsection*{Parameters}

\begin{center}\footnotesize
\begin{tabular}{@{}lp{0.7\textwidth}@{}}
\toprule
Parameter & Meaning \\
\midrule
\code{DIV} & The bit period \code{div} starts with, in cycles; the
register resets to \code{DIV} less one. The tables use 868, which
\code{uart.rs} states is 115200 baud at 100~MHz, as the board runs.
The checked runs use 4. \\
\bottomrule
\end{tabular}
\end{center}

\subsection*{Register map}

Offsets from the port's base. On the board the base is \code{0x3000},
which is also \code{isa::UART\_BASE}. The port selects the word by
bits 4 to 2 of the address.

\dsregs{Uart}

The map is declared as \code{serial}. Bit 31 of \code{txdata} reads
one while the transmit queue is full, and a byte written then is
dropped. A read of \code{rxdata} takes the oldest byte; bit 31 reads
one when there was none, and the low byte then reads zero. \code{ip}'s
\code{txwm} is set while fewer than \code{txcnt} bytes wait to go
out, and its \code{rxwm} while more than \code{rxcnt} wait to be read.
Writes to \code{rxdata} and \code{ip} are ignored, and \code{0x1c},
which no register names, reads zero.

Two resets differ from SiFive's, so that a program that never writes
the controls finds the port running: \code{txen} and \code{rxen}
start set, and \code{div} starts at \code{DIV} less one, 867 for
115200 baud at 100~MHz. \code{ie} starts clear, as SiFive's does, so
\code{irq} stays low until a program sets a watermark's enable. It
started with \code{rxwm} set until issue~\#1137: Linux's driver asks
for its interrupt before its port is registered, and a receive enable
that started set let the loader's bytes leave a request pending in the
interrupt controller, which Linux took before its port was ready
(issue~\#1136).

\dstables{Uart}

\subsection*{Behaviour}

\begin{itemize}
\item A read is taken when \code{bus\_r} has room and a request is on
\code{bus\_ar}, and it is answered with \code{OKAY} in the same cycle.
\item A write is taken when \code{bus\_b} has room and both
\code{bus\_aw} and \code{bus\_w} hold a request. It is answered with
\code{OKAY} in the same cycle, whichever word it names.
\item A byte written to \code{txdata} with the queue not full goes in
at the queue's tail. It also goes into \code{last}, and \code{sent}
counts it.
\item When the line is idle, the queue holds a byte and \code{txen}
is set, the oldest byte leaves the queue. Its frame loads into
\code{shift}, and \code{bits} into ten, or eleven with \code{nstop}.
\item While \code{bits} is not zero the port is busy and \code{tx} is
bit 0 of \code{shift}. Every \code{div} plus one cycles the frame
shifts down and \code{bits} counts down. When idle, \code{tx} is high.
\item The \code{rx} line passes through one register, \code{line},
before the logic.
\item When the port is not receiving, \code{rxen} is set and
\code{line} is low, a frame begins. The port samples \code{line} in
the middle cycle of each of the ten bits.
\item A high where the start bit should be is a false start, and the
frame is dropped. The eight data bits shift into \code{rx\_shift}.
\item A high stop bit lands the byte. If the receive queue holds
sixty-four, the byte is dropped and \code{dropped} counts it; else it goes
in at the tail and \code{received} counts it. A low stop bit lands
nothing.
\item A push and a removal in the same cycle leave either queue's
count unchanged.
\item \code{irq} is \code{ie}'s \code{txwm} and \code{ip}'s
\code{txwm}, or \code{ie}'s \code{rxwm} and \code{ip}'s \code{rxwm}.
\item With \code{rst} high, \code{bits}, \code{tick},
\code{rx\_bits} and \code{rx\_tick} go to zero, both queues empty, and
the controls and \code{div} return to their reset values.
\end{itemize}

\subsection*{Verification}

\begin{itemize}
\item \code{//cpu/vreteno:lockstep\_test} puts the terminal of
\code{term.rs} on the two lines. At the halt it checks \code{sent} and
\code{last} against the bytes the model was given, and that
\code{dropped} is zero. In \code{demo\_program} the port must say
\code{OK} and a newline, receive three bytes, and echo them.
\item \code{//cpu/vreteno:hello\_test},
\code{//cpu/vreteno:hello\_cc\_test} and \code{//cpu/vreteno:board\_test}
check what compiled programs print on \code{tx}. The harnesses run the
line on after the halt, while the transmit queue empties.
\item \code{//cpu/vreteno:board\_test}'s
\code{a\_line\_typed\_back\_to\_back\_waits\_whole\_in\_the\_receive\_queue}
types forty-eight bytes at the line's full speed while the program is
busy, and only then has it read and echo them. Every byte must come
back, which a queue of eight would not allow.
\item \code{//cpu/vreteno:zephyr\_port\_test} checks that the map's
offsets and masks are the ones Zephyr's stock \code{uart\_sifive}
driver uses, that Zephyr's device tree names the port
\code{sifive,uart0} with its clock, and that the console is that port.
\item The build simulates \code{Uart<4>}'s netlist, the entity
\code{uart4}, against the trace of \code{//cpu/vreteno:vreteno} under
nvc and Verilator: \code{//docs:vreteno\_sim\_uart4\_uart4\_tb\_test}
and \code{//docs:vreteno\_vsim\_uart4\_test}. The run also writes
\code{Uart<868>}'s netlist, which no test simulates.
\end{itemize}

\subsection*{Used by}

\code{Board}, as \code{uart} behind the bridge \code{puart}, with
\code{irq} as source 1 of the interrupt controller. The
demonstration \code{//cpu/vreteno:vreteno}, the lockstep test and
\code{run.rs}. The system in \code{soc}, where it is the serial port at
corner 1,1 of the lattice. Zephyr's console on Vreteno, through the
stock \code{uart\_sifive} driver.

\subsection*{Limits}

A byte written while the transmit queue is full is dropped, and a
byte received with the receive queue full is dropped and counted.
\code{sent}, \code{received} and \code{dropped} are eight-bit counters
that are not on the bus. The port checks no address beyond bits 4 to
2, since the bridge sends it only its own range. There is no parity
bit and no flow-control line.
